\documentclass{beamer}
\usetheme{STCE}
\usepackage{beamerfoils}

\usepackage{epstopdf}

\usepackage{ucs} % Unicode - dependency of utf8x inputenc
\usepackage[utf8x]{inputenc} % "utf8x" uses "ucs"-package, better than "UTF8"
\usepackage[T1]{fontenc}
\usepackage{lmodern}

\usepackage[german]{babel}
\usepackage{amsmath, amssymb, amsxtra, amsthm}
\usepackage{epsfig,psfrag, epstopdf}
\usepackage{listings}

\usepackage{tikz}

\usepackage{nicefrac}


\beamertemplatenavigationsymbolsempty % turn off navigation bar


\setbeamertemplate{footline}[frame number]

\MyLogo{
\centering
\includegraphics[width=0.12\textwidth]{./figures/logo.eps}
}

\begin{document}
\title{\centering
\includegraphics[width=0.25\textwidth]{./figures/logo.eps} \\ Numerische Bibliotheken}

\subtitle{Der $\chi^2$ und Kolmogorow-Smirnow Test}
\author{Christian Janßen (302530) \\ Fabian Ohler (280424) }
\date{6. August 2013}
\institute{
LuFG Informatik 12: Software and Tools for Computational Engineering}
\frame[plain]{\titlepage}

%%% Inhalt
\begin{frame}
\frametitle{Inhalt}
	\begin{itemize}
		\item Statistische Grundlagen
%		\begin{itemize}
%			\item test 1.2
%			\item test 1.1
%		\end{itemize}
		\item Kolmogorow-Smirnow Test
		\begin{itemize}
			\item Mathematische Hintergrund
			\item Implementierung in NAG C
			\item Alternative Bibliotheken
		\end{itemize}
		\item $\chi^2$ Test
		\begin{itemize}
			\item Mathematische Hintergrund
			\item Implementierung in NAG C
			\item Alternative Bibliotheken
		\end{itemize}
		\item Live Demo
		\item Zusammenfassung
	\end{itemize}
\end{frame}

%%%Statistische Grundlagen
\begin{frame}
\frametitle{Statistische Grundlagen}
\end{frame}

%%%Kolmogorow-Smirnow Test
\begin{frame}
\frametitle{Kolmogorow-Smirnow Test}
Kolmogorow-Smirnow Test hat zwei Anwendungen:\\
	\begin{itemize}
		\item[1.] Pr\"uft Mengen $X$ und $ Y$ auf gleiche Wahrscheinlichkeitsverteilung \\
		\item[2.] Pr\"uft Menge $X$ auf spezifizierte Wahrscheinlichkeitsverteilung $F(x)$ \\
			\[S(x)=F(x)\]
			\begin{center}$\rightarrow$ Anpassungsg\"utetest\end{center}
	\end{itemize}
\end{frame}
	%%%Mathematischer Hintergrund
\begin{frame}
\frametitle{Kolmogorow-Smirnow Test}
Anwendung des Anpassungsg\"utetests:
	\begin{itemize}
		\item Menge $X$ aufsteigend sortieren, $x_1\le ...\le x_n$
		\item<2-> Bilden der empirischen Verteilungsfunktion $S(x)$	
		\only<2>{ $$S(x)=\frac{1}{n}\sum^{n}_{i=1}I_{x_i\le x}$$}
		\item<3-> Spezifizieren der angenommenen Verteilungsfunktion $F(x)$
		\item<3-> Aufstellen der Nullhypothese $H_0$ und der Alternativhypothese $H_1$
		\only<3>{$$H_0:S(x)=F(x)$$ $$H_1:S(x)\neq F(x)$$}
		\item<4-> Ermitteln der Teststatistik $D_n$
		\only<4>{ $$D_n=\max(d_0,d_1)$$ $$ d_0=\max_x|S(x)-F(x)|,$$ $$d_1=\max_x|S(x_{i-1})-F(x_i)|$$ }
		\item<5->Wenn $D_n>D_\alpha$ wird die Nullhypothese verworfen
	\end{itemize}
\end{frame}
	%%%Beispie
\begin{frame}
\frametitle{Kolmogorow-Smirnow Test}
Ein Beispiel:\\
	\begin{itemize}
		\item Beobachtungen:
			\begin{table}[h]
			\center
			\begin{tabular}{c|c|c|c|c|c|c|c|c|c}
			$x_1$	&$x_2$	&$x_3$	&$x_4$	&$x_5$	&$x_6$	&$x_7$	&$x_8$	&$x_9$	&$x_{10}$\\
			\hline
			1	&1.1	&1.2	&1.6	&1.7	&2.1	&2.1	&2.4	&2.4	&2.5	\\
			\end{tabular}
			\begin{tabular}{c|c|c|c|c|c|c|c|c|c}
			$x_{11}$	&$x_{12}$	&$x_{13}$	&$x_{14}$	&$x_{15}$	&$x_{16}$	&$x_{17}$	&$x_{18}$	&$x_{19}$	&$x_{20}$\\
			\hline
			2.6	&2.6	&2.6	&2.7	&2.8	&3	&3.3	&3.5	&3.8	&4.2\\
			\end{tabular}
			\end{table}
		\item Hypothesen: $$H_{Franz}:S(x)=\Phi(x|3.5;1)$$ $$H_{Paul}:S(x)=\Phi(x|3;0.7)$$
	\end{itemize}
\end{frame}

\begin{frame}
\frametitle{Kolmogorow-Smirnow Test}
Berechnung der Werte:\\
	\begin{itemize}
		\item Werte der Verteilungsfunktionen
			\begin{table}[ht]
			\center
			\begin{tabular}{c|c|c|c|c}
			$i$ 	& $x_i$ 	& $S(x_i)$ 	& $\Phi (x_i|3;0,7)$ 	& $\Phi (x_i|3,5;1)$ 	\\
			\hline
			1	&	1	&	0.05	&	0.002137	&	0.006210	\\
			2	&	1.1	&	0.1	&	0.003321	&	0.008198	\\
			%3	&	1.2	&	0.15	&	0.005064	&	0.010724	\\
			%4	&	1.6	&	0.2	&	0.022750	&	0.028717	\\
			%5	&	1.7	&	0.25	&	0.031645	&	0.035930	\\
			%6	&	2.1	&	0.35	&	0.099271	&	0.080757	\\
			%7	&	2.1	&	0.35	&	0.099271	&	0.080757	\\
			%8	&	2.4	&	0.45	&	0.195683	&	0.135666	\\
			%9	&	2.4	&	0.45	&	0.195683	&	0.135666	\\
			%10	&	2.5	&	0.5	&	0.237525	&	0.158655	\\
			%11	&	2.6	&	0.65	&	0.283855	&	0.184060	\\
			12	&	2.6	&	0.65	&	0.283855	&	0.184060	\\
			%13	&	2.6	&	0.65	&	0.283855	&	0.184060	\\
			%14	&	2.7	&	0.7	&	0.334118	&	0.211855	\\
			15	&	2.8	&	0.75	&	0.387548	&	0.241964	\\
			%16	&	3	&	0.8	&	0.500000	&	0.308538	\\
			%17	&	3.3	&	0.85	&	0.665882	&	0.420740	\\
			%18	&	3.5	&	0.9	&	0.762475	&	0.500000	\\
			%19	&	3.8	&	0.95	&	0.873451	&	0.617911	\\
			20	&	4.2	&	1	&	0.956762	&	0.758036	\\
			\end{tabular}
			\end{table}
		\item
\begin{table}[ht]
\center
\begin{tabular}{c|c|c|c|c}
$i$ 	& $d_{0,Paul}(x_i)$ 	& $d_{1,Paul}(x_i)$ 	& $d_{0,Franz}(x_i)$ 	& $d_{1,Franz}(x_i)$ 	\\
\hline
1	&	0.047863	&	0.002137	&	0.043790	&	0.006210	\\
2	&	0.096679	&	0.046679	&	0.091802	&	0.041802	\\
..	&	...		&	...		&	...		&	...		\\
%3	&	0.144936	&	0.094936	&	0.139276	&	0.089276	\\
%4	&	0.177250	&	0.127250	&	0.171283	&	0.121283	\\
%5	&	0.218355	&	0.168355	&	0.214070	&	0.164070	\\
%6	&	0.250729	&	0.150729	&	0.269243	&	0.169243	\\
%7	&	0.250729	&	0.250729	&	0.269243	&	0.269243	\\
%8	&	0.254317	&	0.154317	&	0.314334	&	0.214334	\\
%9	&	0.254317	&	0.254317	&	0.314334	&	0.314334	\\
10	&	0.262475	&	0.212475	&	0.341345	&	0.291345	\\
11	&	0.366145	&	0.216145	&	0.465940	&	0.315940	\\
12	&{0.9}	0.366145	&	0.366145	&	0.465940	&	0.465940	\\
..	&	...		&	...		&	...		&	...		\\
%13	&\cellcolor[gray]{0.9}	0.366145	&	0.366145	&	0.465940	&	0.465940	\\
%14	&	0.365882	&	0.315882	&	0.488145	&	0.438145	\\
15	&	0.362452	&	0.312452	&[gray]{0.9}	0.508036	&	0.458036	\\
%16	&	0.300000	&	0.250000	&	0.491462	&	0.441462	\\
%17	&	0.184118	&	0.134118	&	0.429260	&	0.379260	\\
%18	&	0.137525	&	0.087525	&	0.400000	&	0.350000	\\
..	&	...		&	...		&	...		&	...		\\
19	&	0.076549	&	0.026549	&	0.332089	&	0.282089	\\
20	&	0.043238	&	0.006762	&	0.241964	&	0.191964	\\
\end{tabular}
\end{table}


	\end{itemize}
\end{frame}
	%%%NAG C
\begin{frame}
\frametitle{Kolmogorow-Smirnow Test}
\end{frame}
	%%%R
\begin{frame}
\frametitle{Kolmogorow-Smirnow Test}
\end{frame}

%%%$\chi^2$ Test
\begin{frame}
\frametitle{$\chi^2$ Test}
\end{frame}

%%%Live Demo
\begin{frame}
\frametitle{Live Demo}
\center
\Huge Live Demo
\end{frame}

%%%Zusammenfassung
\begin{frame}
\frametitle{Zusammenfassung}
\end{frame}


\end{document}
